\chapter{LITERATURE REVIEW}

The rapid growth of digital communication technologies has significantly transformed the way people interact and collaborate. Video conferencing systems have become essential tools in domains such as education, business, healthcare, and remote work environments. This chapter reviews existing technologies and research contributions related to web-based video conferencing systems.

Early video conferencing solutions relied heavily on proprietary software and dedicated hardware, making them expensive and less accessible. Platforms such as Skype introduced internet-based calling, but scalability and real-time performance remained challenges. With advancements in web technologies, modern systems have shifted toward browser-based solutions using protocols like WebRTC (Web Real-Time Communication), which enables peer-to-peer communication without requiring additional plugins.

Recent platforms such as Zoom, Google Meet, and Microsoft Teams have set benchmarks in usability, scalability, and feature richness. These platforms provide functionalities such as real-time audio/video communication, screen sharing, meeting recording, and participant management. Their success highlights the importance of low latency, high reliability, and intuitive user interfaces in modern communication systems.

Research in this field has focused on improving network efficiency, reducing latency, and ensuring secure communication. Techniques such as adaptive bitrate streaming and congestion control algorithms have been widely studied to maintain call quality under varying network conditions. Additionally, security mechanisms like end-to-end encryption have been emphasized to protect user data and privacy.

With the emergence of modern web frameworks like Next.js and real-time communication services such as GetStream, developers can now build scalable and efficient applications more easily. Authentication systems like Clerk further enhance security by providing robust identity management solutions.

The proposed system builds upon these advancements by integrating modern web technologies to develop a responsive and secure video conferencing application. It combines real-time communication capabilities with user-friendly design and scalable architecture, addressing the limitations of earlier systems while aligning with current industry standards.

\section{Video Conferencing Platforms}

Popular platforms such as Zoom and Google Meet provide seamless communication using cloud infrastructure. These platforms focus on scalability, reliability, and user experience.

\section{WebRTC Technology}

WebRTC (Web Real-Time Communication) is a key technology that enables peer-to-peer communication directly between browsers. It supports audio, video, and data sharing without requiring external plugins.

\section{Authentication Systems}

Modern applications use authentication services such as OAuth, JWT, and third-party providers to ensure secure access. Clerk provides a robust authentication solution with support for social logins.

\section{Real-Time Communication Systems}

Real-time systems require low latency, efficient data transfer, and synchronization across multiple users. Technologies such as WebSockets and streaming APIs are commonly used.

\section{Challenges}

\begin{itemize}
\item Network latency and bandwidth limitations
\item Security concerns in data transmission
\item Scalability for large number of users
\end{itemize}

\section{Summary}

The literature review highlights the importance of scalable architecture, secure authentication, and efficient real-time communication in building video conferencing systems.